\section{Research Agenda and Limitations}
\label{sec:agenda}

\subsection{Action Semantics and Causal Identification}

The first priority is to define actions before modeling their effects. Medical datasets contain treatment orders, administrations, dose changes, procedures, monitoring choices, clinician recommendations, and workflow events. These variables differ in executability, timing, adherence, and causal meaning. Future work should report the action set, admissible range, decision time, provenance, and empirical support, and should distinguish an observed care event from an action that can be set in simulation. This simple discipline would resolve many Level~2/Level~3 boundary disputes and expose when a proposed intervention lies outside the data.

Action-aware simulation should also be paired with negative controls. Holding state fixed while varying treatment, masking action inputs, comparing with time-only models, and testing known dose-response or physical relationships can reveal whether the model uses the intervention as intended. These tests do not establish causality, but they prevent action tokens from serving as decorative conditioning. In EHR systems, separating the patient transition from the clinician or observation policy is especially important. Joint models can be appropriate, but the modeled process should be named.

The next step is to connect flexible world models with causal identification. Counterfactual sequence models, trial emulation, instrumental or quasi-experimental designs, mechanistic constraints, and sensitivity analyses should not remain peripheral to medical-world-model development. A productive integration would retain rich multimodal state and generative uncertainty while making the estimand, support assumptions, censoring, and time-varying confounding explicit. The aim is not to force every simulator to claim treatment effects. It is to make clear when such a claim is intended and what evidence it requires.

\subsection{Decision-Aligned Evaluation, Uncertainty, and External Validation}

Evaluation should follow the decision horizon. One-step image or event prediction may be sufficient for interpolation but not for long-range planning. A treatment simulator should be evaluated on action sensitivity and clinically relevant response, not only perceptual quality. A planner should be assessed on how transition error changes selected actions, not only on simulator reward. Each capability level therefore needs a distinct evaluation contract rather than a shared leaderboard.

Uncertainty must constrain use. Generative diversity is not enough; systems should estimate when state, action, or horizon lies outside support. Calibration, coverage, conformal or ensemble methods, model disagreement, and mechanistic sensitivity can all contribute, but the output should affect the decision loop. A planner that recognizes uncertainty yet still optimizes aggressively outside the training distribution has not solved the problem. Abstention, conservative fallback policies, and uncertainty-aware action constraints should be evaluated explicitly.

External validation is particularly important because clinical dynamics are shaped by population, site, scanner, coding, workflow, and treatment policy. A model trained at one institution may fail because the patient distribution changes or because the observation and care processes differ. External cohorts should therefore report which component shifted. For action-aware systems, differences in treatment practice can be informative rather than merely inconvenient: they test whether the model learned patient response or local policy. Prospective silent deployment and workflow studies can provide an intermediate step between retrospective evaluation and interventional trials.

The capability-specific requirements above are summarized as a validation ladder in Appendix~\ref{app:coding}. Its rungs are evidence types rather than a universal score: the appropriate evidence depends on the intended use, and one rung does not substitute for causal identification or action support at another.

\subsection{Data Coverage, Generalization, and Hybrid Mechanistic Models}

Current work concentrates in settings with repeated measurements, large archives, or simulatable environments. Rare diseases, multimorbidity, long-term outcomes, underrepresented populations, and cross-system trajectories remain less visible. Scaling the same architecture on more records will not automatically solve this gap because actions and outcomes are often sparsely observed. Data collection should be organized around decisions and longitudinal follow-up, with explicit documentation of treatment timing, adherence, censoring, and missingness.

Multimodal integration should be selective. Imaging, EHR, genomics, physiology, pathology, behavior, and environment can all contribute to patient state, but broad fusion is useful only when each modality supports the transition or outcome being modeled. A narrow, validated simulator can be more scientifically informative than a universal patient representation whose action pathway is unclear. Missing-modality behavior and asynchronous measurement should be treated as part of the model rather than hidden by imputation.

Hybrid mechanistic models are a promising response to sparse interventions and distribution shift. Known anatomy, conservation laws, dose-response relationships, pharmacokinetics, tissue mechanics, and safety constraints can restrict implausible rollouts. Learned components can estimate residual dynamics, infer patient parameters, or process complex observations. The key challenge is calibration: structural assumptions and parameter uncertainty must be reported, and the model should be tested where learned and mechanistic components disagree. Hybrid does not mean valid by construction.

\subsection{Benchmarks and Reporting Standards}

Benchmarks should be stratified by capability rather than mixing static prediction, natural-history forecasting, action-conditioned simulation, counterfactual estimation, and planning. For Level~2, tasks should distinguish interpolation, one-step prediction, and long-horizon rollout. For Level~3, they should define executable actions, hold pre-action state fixed, and include no-action and action-shuffle controls. For Level~4, they should state the estimand and provide trial, quasi-experimental, semi-synthetic, or mechanistic anchors with overlap and sensitivity analyses. For Level~5, they should evaluate policy robustness, reward sensitivity, support violations, and the relationship between model error and decision error.

A compact reporting card can make claims comparable without prescribing one model family. It should state the modeled state and time scale; action variables, timing, and provenance; transition mechanism; outcome and clinical horizon; uncertainty; data and action support; validation setting; and intended use. Authors should identify the capability level they implement and separately state whether causal or clinical validity has been evaluated. Reviewers can then judge a narrow claim on its own evidence rather than rewarding the most ambitious vocabulary.

Adjacent evidence records clarify what a benchmark cannot establish. A multicenter longitudinal cohort can characterize prognostic trajectories and outcome associations without providing ground truth for an intervention simulator \citep{Li2021TrajectoriesPerioperativeSerum}. Conversely, a digital-twin framework can optimize secure data management without implementing patient-state dynamics \citep{Alotaibi2026CostOptimizedTwin}. Such studies remain relevant to data provenance and evaluation design, but they should not be counted as capability implementations. Reporting the paper role alongside the benchmark task prevents infrastructure, cohort evidence, and simulated dynamics from being conflated.

Open code and data improve auditability but should not be treated as proxies for scientific quality. Clinical data may be restricted for legitimate reasons. When data cannot be released, cohort construction, event vocabularies, preprocessing, censoring, splits, and evaluation protocols become more important. Paper-linked code, model weights, synthetic examples, and executable evaluation scripts can still support reproducibility. Resource status should be timestamped because repositories and releases change.

\subsection{Review Limitations}

This survey is an ML-facing capability and evidence review, not a disease-specific clinical systematic review, regulatory assessment, formal risk-of-bias analysis, or meta-analysis. Terminology creates retrieval bias: papers using ``world model'' are easier to find than methods with equivalent task contracts under different names. We mitigate this by including adjacent digital-twin, trajectory, causal, and model-based decision literatures, but the boundary cannot be exhaustive.

The capability taxonomy is adapted from prior medical-world-model frameworks, and its categories remain analytical choices rather than a settled community standard \citep{Qazi2025BeyondGenerativeAI,Saeed2026MedicalWorldModel,Liu2026MedicalWorldModelsReview,Wang2026TowardsBiomedicalWorldModels}. The non-cumulative interpretation, clinical-action boundary, and SATO-V evidence contract are designed to make those choices explicit. Some papers remain reasonable boundary cases, particularly when actions are partially implemented or causal assumptions are implicit. The detailed coding protocol preserves uncertainty rather than forcing every study into an overprecise label.

The literature is also moving quickly. Preprints, publication status, repositories, and validation evidence can change. The corpus therefore supports a current evidence synthesis rather than a final census of the field. Architecture is discussed qualitatively because a normalized codebook has not yet been validated across the full literature; no prevalence claim is made from free-text descriptions. Finally, public evidence may understate private implementation or clinical testing. Throughout the review, ``not publicly evidenced'' should not be read as proof of absence.

These limitations do not remove the main conclusion. They reinforce the need to separate task capability from evidential strength and to make boundary judgments inspectable. Formal submission should be preceded by author-level verification of every included record, refreshed metadata and resources, and domain review of the strongest clinical claims.
